\chapter{Risk Vectors, Weak Elements, and Planning Metrics}

The damping-operator hierarchy changes how weak elements and planning actions
should be defined. A weak bus, weak line, or useful converter placement is not
an intrinsic property of a graph alone. It is a property relative to a dynamic
margin, an operator class, an action family, and a validation gate. This
chapter turns the spectral derivatives into a planning language.

\section{Margins Must Name the Mechanism}

The most common small-signal margin is the critical damping ratio,
\begin{equation}
  S_\zeta=\log\frac{\zetaReq}{\zeta_{\rm crit}}.
  \label{eq:imp_zeta_margin}
\end{equation}
The case satisfies the requirement when \(S_\zeta\le0\). This margin is natural
for oscillatory eigenvalues, but it does not measure non-normal amplification,
input-output gain, or protection-weighted response.

For a state matrix \(A\), selected inputs \(B\), selected outputs \(C\), and
weights \(\Gamma_u,\Gamma_y\), a protected resolvent gain can be defined as
\begin{equation}
  g_{\rm prot}=
  \sup_{\omega\in\Omega}
  \sigma_1\left(\Gamma_yC(j\omega I-A)^{-1}B\Gamma_u\right).
  \label{eq:imp_protected_gain}
\end{equation}
The corresponding log margin is
\begin{equation}
  S_g=\log g_{\rm prot}.
  \label{eq:imp_Sg}
\end{equation}
The two margins can rank actions differently. That disagreement is not an
error; it means the mechanisms being measured are different.

\section{Action Coordinates}

Let a planning action be parameterized as
\begin{equation}
  \rho=
  \begin{bmatrix}
  \rho_L\\
  \rho_D\\
  \rho_M
  \end{bmatrix},
  \label{eq:action_vector}
\end{equation}
where \(\rho_L\) changes line weights or topology, \(\rho_D\) changes damping
or primary-control gain, and \(\rho_M\) changes physical or virtual inertia.
For a simple pole \(s_\star\), the total first-order movement is
\begin{equation}
  ds_\star
  =
  \nabla_Ls_\star^T d\rho_L
  +
  \nabla_Ds_\star^T d\rho_D
  +
  \nabla_Ms_\star^T d\rho_M.
  \label{eq:mixed_gradient}
\end{equation}
The three gradients are computed from the exact QEP sensitivity
\eqref{eq:fixed_full_qep_sens} or, after condensation, from the NEP
sensitivity \eqref{eq:self_nep_sens}. The key point is that the gradients are
not interchangeable. \(L\)-moves are stiffness moves, \(D\)-moves are
dissipation moves, and \(M\)-moves are inertia and normalization moves.

\section{Finite-Action Coordinates}

The action coordinates above are local, but the rational self-energy chapter
gives a finite-action object. For any low-rank action
\[
  T_u(s)=T(s)+U(s)K_u(s)V(s)^*,
\]
the action interaction matrix is
\begin{equation}
  \mathcal G_u(s)=V(s)^*T(s)^{-1}U(s),
  \label{eq:imp_action_green_matrix}
\end{equation}
and the protected roots created or shifted by the action satisfy
\begin{equation}
  \det\!\left[I+K_u(s)\mathcal G_u(s)\right]=0.
  \label{eq:imp_action_root_condition}
\end{equation}
This formula changes how a planning action should be reported. A derivative
\(ds_\star/du\) is the tangent of \eqref{eq:imp_action_root_condition}; it is
not the full action law. For a finite line increment, the scalar interaction is
the dynamic effective resistance
\(\mathcal R_e(s)=a_e^TT(s)^{-1}a_e\). For finite damping, inertia, and
controller-retuning actions, the relevant object is the corresponding block of
\(\mathcal G_u(s)\).

A local risk vector should therefore be interpreted as the first-order
linearization of three finite objects:
\begin{equation}
  \hbox{action determinant}
  \quad
  \hbox{protected margin boundary}
  \quad
  \hbox{structured fragility radius}.
  \label{eq:imp_finite_action_objects}
\end{equation}
This interpretation prevents a candidate from being accepted only because its
infinitesimal derivative is favorable at the base point. The finite action must
remain inside the protected margin and outside the structured pseudospectral
boundary over the declared budget.

\section{Risk Vector for a Candidate Action}

For any candidate action \(a\), define a local risk vector after normalizing the
action coordinate and each margin. A gain component must be dimensionless, e.g.
\(S_g=\log(g_{\rm prot}/g_{\rm req})\), not the logarithm of a dimensional
quantity. With a declared action scale \(\bar a\), derivatives below are taken
with respect to \(\hat a=a/\bar a\):
\begin{equation}
  r(a)=
  \begin{bmatrix}
  r_\sigma(a)\\
  r_\zeta(a)\\
  r_g(a)\\
  r_h(a)\\
  r_f(a)
  \end{bmatrix}.
  \label{eq:risk_vector}
\end{equation}
The components are:
\begin{align}
  r_\sigma(a)
  &=
  \left[
  \frac{d\Real(s_{\rm crit})}{da}
  \right]_+,
  \label{eq:risk_sigma}\\
  r_\zeta(a)
  &=
  \left[
  \frac{dS_\zeta}{da}
  \right]_+,
  \label{eq:risk_zeta}\\
  r_g(a)
  &=
  \left[
  \frac{dS_g}{da}
  \right]_+,
  \label{eq:risk_gain}\\
  r_h(a)
  &=
  \mathbf 1\left\{
  \frac{d\nu_k}{da}>0,\ 
  \frac{d\zeta_k}{da}<0
  \right\},
  \label{eq:risk_hidden}\\
  r_f(a)
  &=
  \hbox{normalized feasibility penalty}.
  \label{eq:risk_feasibility}
\end{align}
Here \([x]_+=\max(x,0)\). The first entry detects rightward pole movement. The
second detects damping-ratio margin loss. The third detects protected
input-output amplification. The fourth detects the hidden-margin conflict from
Chapter~4. The fifth records voltage, flow, current, protection, or dispatch
constraints.

All entries in \eqref{eq:risk_vector} are dimensionless local rates after the
chosen normalization. Without that convention, the angle between two risk
vectors can change simply because the same action is expressed in MW, per-unit,
or percent penetration.

A scalar screening score may be formed as
\begin{equation}
  \mathcal R(a)=w^Tr(a),\qquad w\ge0,
  \label{eq:scalar_risk_score}
\end{equation}
but the vector should be reported before scalarization. Otherwise, a strong
graph improvement can hide a damping loss.

\section{Bus Risk Vector}

For bus \(i\), define a fixed-\(D\) spectral bus vector
\begin{equation}
  r_i^{\rm bus}=
  \begin{bmatrix}
  \mathcal D_i\\
  \mathcal M_i\\
  \Pi_i\\
  H_i
  \end{bmatrix},
  \label{eq:bus_risk_vector}
\end{equation}
where
\begin{equation}
  \mathcal D_i=
  \sum_{k\in\mathcal K}
  \omega_k
  \frac{\phi_k(i)^2}{2\sqrt{\nu_k}}
  \label{eq:bus_D_score}
\end{equation}
is damping placement leverage,
\begin{equation}
  \mathcal M_i=
  \sum_{k\in\mathcal K}
  \omega_k\zeta_k\phi_k(i)^2
  \label{eq:bus_M_score}
\end{equation}
is inertia-dilution exposure under fixed \(D\),
\begin{equation}
  \Pi_i=
  \sum_{k\in\mathcal K}\omega_k\phi_k(i)^2
  \label{eq:bus_participation_score}
\end{equation}
is modal participation exposure, and \(H_i\) records whether incident graph
actions produce hidden damping conflicts. A high \(\mathcal D_i\) says
``put damping here.'' A high \(\mathcal M_i\) says ``do not put inertia here
without co-designed damping if the objective is damping ratio.'' A high
\(\Pi_i\) says the bus is dynamically visible.

\section{Line Risk Vector}

For a line \(e=(i,j)\), define
\begin{equation}
  r_e^{\rm line}=
  \begin{bmatrix}
  G_e^{(2)}\\
  G_e^{(K)}\\
  Z_e\\
  H_e\\
  F_e
  \end{bmatrix}.
  \label{eq:line_risk_vector}
\end{equation}
The algebraic-connectivity gain is
\begin{equation}
  G_e^{(2)}=(\phi_2(i)-\phi_2(j))^2.
  \label{eq:line_fiedler_gain}
\end{equation}
The effective-resistance gain is
\begin{equation}
  G_e^{(K)}
  =
  a_e^T(L^+)^2a_e.
  \label{eq:line_kirchhoff_gain}
\end{equation}
The damping-ratio risk for a monitored mode \(k\) is
\begin{equation}
  Z_e^{(k)}
  =
  -\frac{\partial\zeta_k}{\partial b_e},
  \qquad
  Z_e=\max_{k\in\mathcal K} Z_e^{(k)}.
  \label{eq:line_zeta_risk}
\end{equation}
The hidden-margin flag is
\begin{equation}
  H_e=\mathbf 1\left\{
  G_e^{(2)}>0,\ Z_e>0
  \right\},
  \label{eq:line_hidden_flag}
\end{equation}
and \(F_e\) is a feasibility penalty for thermal, voltage, protection, or
right-of-way constraints. A high \(G_e^{(2)}\) or \(G_e^{(K)}\) is not enough:
the line is attractive only if \(Z_e\) and \(F_e\) are acceptable.

\section{Weak Nodes Are Relative Objects}

A bus is weak only after the following items are specified:
\begin{enumerate}
\item the base operating point and dynamic model;
\item the operator class being used: uniform \(D=\gamma M\), commuting fixed
\(D\), non-commuting fixed \(D\), rational \(\Pi(s)\), or damping-channel
\(\Sigma(s)\);
\item the margin: damping ratio, decay rate, protected resolvent gain, or
another registered metric;
\item the action family: line change, damping addition, inertia change,
converter placement, controller retuning, or protection setting;
\item the accept/reject gate that turns a screen into evidence.
\end{enumerate}

Without those items, a weak-node list is not reproducible. A bus may be weak
for damping-ratio improvement under controller retuning and irrelevant for
protected-gain reduction under line reinforcement.

\section{Improving the Graph Without Adding Inertia}

When the objective is to improve graph robustness without adding inertia, the
candidate action should be selected from graph derivatives and then passed
through damping gates. A conservative procedure is:
\begin{enumerate}
\item rank candidate lines by \(G_e^{(2)}\) if the target is algebraic
connectivity;
\item rank candidate lines by \(G_e^{(K)}\) if the target is total effective
resistance or broad slow-mode coherence;
\item remove candidates with unacceptable \(F_e\);
\item remove or pair candidates with damping support if \(H_e=1\);
\item verify the remaining candidates with QEP or NEP sensitivity.
\end{enumerate}

This gives a mathematically defensible answer to the question ``how do we
improve \(L\) without adding inertia?'' Improve \(L\) where it moves the slow
graph modes most, but do not confuse that with damping unless the damping gate
passes.

\section{Substitutability of Actions}

Two actions are substitutable only if they improve the same margin through
compatible mechanisms. Adding physical damping, changing virtual damping,
placing a grid-forming converter, retuning a PLL, reinforcing a line, or
placing inertia can all move the critical eigenvalue, but they do not generally
act through the same operator term.

Let \(a\) and \(b\) be candidate actions. A local substitutability test is
\begin{equation}
  \cos\angle(r(a),r(b))
  =
  \frac{r(a)^Tr(b)}{\norm{r(a)}\,\norm{r(b)}}.
  \label{eq:risk_vector_angle}
\end{equation}
High alignment means the actions stress or relieve the same risks. Low
alignment means they are not substitutes even if both improve one scalar
metric. If one action improves \(S_\zeta\) by adding damping and another
improves \(S_g\) by reducing resolvent amplification, their risk vectors will
not generally align.

\section{Theory-to-Validation Protocol}

The theoretical hierarchy implies a protocol for future empirical work.

\begin{enumerate}
\item Start with the proportional model only as a baseline and state which
margin it certifies.
\item Measure non-commutation through \(\Gamma=Q^T\Dt Q\) before trusting
graph-mode damping claims.
\item If controller states are eliminated, construct or approximate
\(\Pi(s)\), then inspect its poles, residues, channel decomposition, and any
damping-channel \(\Sigma(s)\) over the frequency band of interest.
\item Use eigenvalue sensitivity to decompose action effects into retained and
condensed contributions.
\item Validate screens against full eigenvalue recomputation before using them
as planning rules.
\end{enumerate}

This protocol prevents a visually plausible graph or time-domain result from
being mistaken for a damping certificate.

\section{Implications for Converter-Rich Systems}

Converter-rich systems make the risk vector more important because converter
controls can supply both damping and destabilizing feedback. A grid-following
PLL, a grid-forming droop loop, a current controller, and a plant controller do
not merely add scalar damping. They add states. If those states are retained,
they appear in the full matrix \(A\). If they are condensed, they appear in
\(\Pi(s)\), or in its declared damping-channel specialization \(\Sigma(s)\).

The planning implication is direct. A reduced model that removes converter
states must show where their spectral effect went. If the answer is a constant
coefficient, the model is making a static approximation. If the answer is a
rational self-energy, the model has preserved the mechanism in the reduced
operator. The practical value of the theory is to make that distinction visible
before numerical ranking begins.
